\documentclass[11pt,a4paper]{article}
\usepackage[utf8]{inputenc}
\usepackage[T1]{fontenc}
\usepackage{hyperref}
\usepackage{graphicx}
\usepackage{listings}
\usepackage{xcolor}
\usepackage{booktabs}
\usepackage{geometry}

\geometry{margin=1in}

\definecolor{codegreen}{rgb}{0,0.6,0}
\definecolor{codegray}{rgb}{0.5,0.5,0.5}
\definecolor{codepurple}{rgb}{0.58,0,0.82}
\definecolor{backcolour}{rgb}{0.95,0.95,0.92}

\lstset{
    backgroundcolor=\color{backcolour},   
    commentstyle=\color{codegreen},
    keywordstyle=\color{magenta},
    numberstyle=\tiny\color{codegray},
    stringstyle=\color{codepurple},
    basicstyle=\ttfamily\footnotesize,
    breakatwhitespace=false,         
    breaklines=true,                 
    captionpos=b,                    
    keepspaces=true,                 
    numbers=left,                    
    numbersep=5pt,                  
    showspaces=false,                
    showstringspaces=false,
    showtabs=false,                  
    tabsize=2
}

\title{LoL Knowledge Graph: From Pipeline Construction to Knowledge Base Expansion}
\author{Knowledge Engineering Lab Report}
\date{\today}

\begin{document}

\maketitle

\section{Introduction}
This report details the construction and expansion of a private Knowledge Base (KB) centered on the \textit{League of Legends} (LoL) universe. The project evolved from a multi-source data extraction pipeline into a large-scale Knowledge Graph aligned with Wikidata, specifically designed to support Knowledge Graph Embedding (KGE) tasks.

\section{Phase 1: Knowledge Base Construction}
The initial phase focused on building a robust local ontology and populating it through automated extraction.

\subsection{Multi-Source Extraction Pipeline}
We implemented a modular pipeline in Python to fetch data from two primary official sources:
\begin{itemize}
    \item \textbf{Official LoL Wiki (MediaWiki API):} Used to extract high-level champion metadata, including roles, official positions, lore, and store costs.
    \item \textbf{Riot Data Dragon (REST API):} Used to fetch granular game data, including base statistics, complete spell tooltips, and skin information for all 174 champions.
\end{itemize}

\subsection{Ontology Design}
The ontology (\texttt{lol\_ontology\_v3.ttl}) was designed using RDF/Turtle syntax. Key features include:
\begin{itemize}
    \item \textbf{Class Hierarchy:} Classes for \texttt{lol:Champion}, \texttt{lol:Spell}, \texttt{lol:Region}, and \texttt{lol:Position}.
    \item \textbf{Advanced Modeling:} Use of \texttt{owl:SymmetricProperty} for relationships like \texttt{isAllyOf} and \texttt{isEnemyOf}, and \texttt{owl:inverseOf} for property mapping.
    \item \textbf{Data Enrichment:} Inclusion of CC effects (Stun, Slow, etc.) and gameplay mechanics extracted via text analysis.
\end{itemize}

\section{Phase 2: Alignment and Entity Linking}
Following Lab 4 requirements, we aligned our private entities with the global \textbf{Wikidata} knowledge base.

\subsection{Entity Linking Strategy}
A dedicated script (\texttt{align\_wikidata.py}) was developed to link private champion URIs to Wikidata URIs (\texttt{wd:Q...}). 
\begin{itemize}
    \item \textbf{Method:} Combination of SPARQL bulk queries (targeting entities present in the LoL narrative universe) and Wikidata Search API fallbacks.
    \item \textbf{Results:} Successful alignment of \textbf{170 out of 174} champions with high confidence.
    \item \textbf{Deliverable:} An \texttt{alignment\_mapping.csv} table recording the mapping and confidence scores.
\end{itemize}

\subsection{Predicate Alignment}
We mapped private predicates to Wikidata properties to ensure semantic interoperability:
\begin{itemize}
    \item \texttt{lol:championName} $\rightarrow$ \texttt{rdfs:label}
    \item \texttt{lol:isFromRegion} $\rightarrow$ \texttt{wdt:P1080} (Narrative universe)
    \item \texttt{lol:playsPosition} $\rightarrow$ \texttt{wdt:P413} (Position played)
\end{itemize}

\section{Phase 3: Knowledge Base Expansion}
To support KGE models, the KB was expanded using an anchored expansion strategy from Wikidata.

\subsection{1-Hop and 2-Hop Expansion}
We executed a recursive expansion script (\texttt{expand\_kb\_2hop.py}):
\begin{enumerate}
    \item \textbf{1-Hop:} Fetching all triplets where our 170 champions were either subject or object in Wikidata.
    \item \textbf{2-Hop:} Expanding the objects found in the first hop (e.g., fetching facts about LoL voice actors, game designers, and related media).
\end{enumerate}

\subsection{Final Cleaning and Integration}
The final Knowledge Base (\texttt{final\_kb.nt}) was generated by merging the private LoL triplets with the expanded Wikidata triplets, removing duplicates, and normalizing the output into N-Triples format.

\section{Final Statistics}
The final Knowledge Base meets and exceeds the requirements for meaningful embedding learning:
\begin{table}[h]
\centering
\begin{tabular}{lr}
\toprule
Metric & Value \\
\midrule
Total Triplets & 48,062 \\
Total Entities & 23,911 \\
Total Relations & 12,283 \\
Alignment Confidence & $\geq 0.70$ \\
\bottomrule
\end{tabular}
\caption{Final Knowledge Base Statistics}
\end{table}

\section{Conclusion}
The project successfully transitioned from a specialized game dataset to a large-scale, aligned Knowledge Graph. With over 48,000 unique triplets, the KB is now ready for training state-of-the-art embedding models like TransE, DistMult, and ComplEx in the subsequent lab sessions.

\end{document}
